%=====================================================================
\section{The DBHE subsurface model}
\label{sec:dbhe}

The Deep Borehole Heat Exchanger is a closed coaxial loop built inside a
repurposed well. Water is injected at the surface down the \emph{annulus}
(between the production casing and a central tubing string), is warmed by the
surrounding rock as it descends, turns at the bottom, and returns up the
\emph{tubing bore}. Nothing is produced from the formation; only heat crosses
the borehole wall.

Two sub-models are coupled. The \emph{formation} model is transient and lives on
a time scale of months to decades: the rock around the well cools progressively
as heat is withdrawn. The \emph{fluid} model is quasi-steady and lives on a time
scale of minutes: the water traverses the well and leaves long before the rock
notices. The whole subsurface exposes exactly two numbers to the rest of the
plant --- the produced temperature $T_{\mathrm{ret}}$ and the wall heat flux
that drives the formation forward one step.

%---------------------------------------------------------------------
\subsection{Formation: radial conduction and its two hypotheses}
\label{sec:dbhe:formation}

Around each axial cell the rock is treated as a homogeneous, radially symmetric
conducting medium with no source:
\begin{equation}
\frac{1}{r}\frac{\partial}{\partial r}\!\left(r\,\frac{\partial\theta}{\partial r}\right)
=\frac{1}{\alpha_e}\frac{\partial\theta}{\partial t}, \qquad
r\in[r_{ce},\,R_{\mathrm{far}}],
\label{eq:dbhe:radial}
\end{equation}
where $\theta(r,z,t)=T-T_\infty(z)$ is the temperature \emph{excess} over the
undisturbed geothermal profile
\begin{equation}
T_\infty(z)=T_{\mathrm{surf}}+G_g z ,
\label{eq:dbhe:Tinf}
\end{equation}
with $T_{\mathrm{surf}}=\SI{20}{\celsius}$ and $G_g=\SI{0.031}{\kelvin\per\metre}$
in the baseline configuration --- so the undisturbed bottom-hole temperature at
the demonstration depth $L=\SI{4500}{\metre}$ is \SI{159.5}{\celsius}. Working
in $\theta$ rather than $T$ is not cosmetic: it makes the initial condition
identically zero and the far-field condition homogeneous, which is what allows
the transform of Section~\ref{sec:dbhe:gitt} to be linear and history-free.

The inner boundary at the borehole wall $r=r_{ce}=\SI{0.10795}{\metre}$ (the
\SI{8.5}{inch} hole radius) carries a prescribed heat flux; the outer boundary
carries a prescribed temperature. Two hypotheses are buried here and both
deserve to be stated with their failure modes.

\paragraph{Hypothesis 1: axial conduction in the rock is negligible.}
Equation~\eqref{eq:dbhe:radial} has no $\partial^2\theta/\partial z^2$ term, so
the $N_Z$ axial cells are thermally independent of one another: each one solves
its own one-dimensional radial problem and they communicate only through the
fluid. The justification is a ratio of driving gradients. With $N_Z=50$ cells
over \SI{4500}{\metre} the cell height is \SI{90}{\metre}, and the undisturbed
temperature difference between neighbouring cells is only
$G_g\times\SI{90}{\metre}=\SI{2.79}{\kelvin}$. The radial problem, by contrast,
is driven by tens of kelvin developed over a diffusion length
$\delta\sim\sqrt{\alpha_e t}$ of a few metres. The axial conductive flux
$k_e\,\Delta T_z/\Delta z$ is therefore smaller than the radial flux
$k_e\,\Delta T_r/\delta$ by two to three orders of magnitude.

\emph{Where it fails.} The argument assumes \emph{smooth} axial variation. It
breaks wherever the axial field develops a genuine kink --- at the very top,
where cold injected water meets warm near-surface rock over a short interval,
and at the bottom turn --- and it breaks if the well is refined to cells of order
the radial diffusion length (metres), at which point the neglected term is no
longer small relative to the retained one. At $N_Z=50$ this is comfortable; the
hypothesis should be re-examined before pushing $N_Z$ past a few hundred.

\paragraph{Hypothesis 2: the domain is truncated at a finite radius
$R_{\mathrm{far}}$, held at the undisturbed temperature.}
The rock is unbounded; the model is not. The code truncates at
\begin{equation}
R_{\mathrm{far}}=\max\!\left(\SI{200}{\metre},\;
3.5\sqrt{4\alpha_e\,t_{\mathrm{total}}}\right),
\label{eq:dbhe:Rfar}
\end{equation}
i.e.\ at several diffusion lengths of the \emph{whole} simulated horizon. For
the baseline $\alpha_e=\SI{1.05e-6}{\metre\squared\per\second}$ and
$t_{\mathrm{total}}=\SI{20}{yr}$ the diffusion length is
$\sqrt{4\alpha_e t_{\mathrm{total}}}=\SI{51.5}{\metre}$ and
$3.5\times\SI{51.5}{\metre}=\SI{180}{\metre}$, so it is the
\SI{200}{\metre} \emph{floor} that is active and $R_{\mathrm{far}}=\SI{200}{\metre}$.

The outer condition imposed there is $\theta(R_{\mathrm{far}},t)=0$: a Dirichlet
condition pinning the far field to the undisturbed geothermal temperature. This
is worth emphasising because it is easy to write ``adiabatic outer boundary''
by reflex, and the two choices are not interchangeable. Under a zero-flux outer
boundary a continuously extracting well has \emph{no} steady state --- the
enclosed rock simply cools without bound --- whereas under the Dirichlet
condition the wall-to-far-field problem possesses the finite steady resistance
\begin{equation}
R_{ss}=\frac{\ln\!\left(R_{\mathrm{far}}/r_{ce}\right)}{2\pi k_e},
\label{eq:dbhe:Rss}
\end{equation}
which the tail correction of Section~\ref{sec:dbhe:tail} depends on entirely.
The code's characteristic equation (below) imposes the Dirichlet form, and
Eq.~\eqref{eq:dbhe:Rss} is used consistently with it. With
$k_e=\SI{2.2}{\watt\per\metre\per\kelvin}$,
$R_{ss}=\SI{0.544}{\kelvin\metre\per\watt}$.

\emph{What it costs.} The truncation turns a slowly-diverging unbounded problem
into a bounded one that eventually reaches steady state. Provided the
disturbance never actually reaches $R_{\mathrm{far}}$, the two are
indistinguishable, and Eq.~\eqref{eq:dbhe:Rfar} is precisely the condition
guaranteeing that. Extend the horizon or raise the flux without re-deriving
$R_{\mathrm{far}}$ and the model starts drawing heat from an artificial infinite
reservoir at the boundary, \emph{over}-predicting sustainable output. The
safeguard is that $R_{\mathrm{far}}$ is computed from $t_{\mathrm{total}}$.

%---------------------------------------------------------------------
\subsection{The generalized integral transform}
\label{sec:dbhe:gitt}

\subsubsection{Eigenvalue problem}
Separating Eq.~\eqref{eq:dbhe:radial} gives the Sturm--Liouville problem
\begin{equation}
\frac{1}{r}\frac{\mathrm{d}}{\mathrm{d}r}\!\left(r\frac{\mathrm{d}\psi}{\mathrm{d}r}\right)
+\beta^{2}\psi=0, \qquad \left.\frac{\mathrm{d}\psi}{\mathrm{d}r}\right|_{r_{ce}}=0, \qquad
\psi(R_{\mathrm{far}})=0 ,
\label{eq:dbhe:SL}
\end{equation}
whose solutions are combinations of the zeroth-order Bessel functions $J_0,Y_0$.
The code builds the eigenfunction so that the inner Neumann condition is
satisfied identically:
\begin{equation}
\psi_n(r)=Y_1(\beta_n r_{ce})\,J_0(\beta_n r)-J_1(\beta_n r_{ce})\,Y_0(\beta_n r),
\label{eq:dbhe:psi}
\end{equation}
since $\mathrm{d}J_0(\beta r)/\mathrm{d}r=-\beta J_1(\beta r)$ makes
$\psi_n'(r_{ce})\propto Y_1 J_1-J_1 Y_1=0$ automatically. The eigenvalues are
then the roots of the remaining (outer) condition,
\begin{equation}
\psi(\beta,R_{\mathrm{far}}) =Y_1(\beta r_{ce})\,J_0(\beta R_{\mathrm{far}}) -J_1(\beta
r_{ce})\,Y_0(\beta R_{\mathrm{far}})=0 .
\label{eq:dbhe:chareq}
\end{equation}
The roots are found by brute force --- a scan of $\num{800000}$ points over
$\beta\in(0,200]\,\si{\per\metre}$ for sign changes, each bracket refined by
Brent's method to a tolerance of $10^{-14}$ --- and the first $N=60$ are kept.
Brute force is the right choice: the roots are computed once and reused for all
$\num{29220}$ time steps, so robustness matters far more than speed. At the
baseline geometry the retained band runs from
$\beta_1=\SI{0.0120}{\per\metre}$ to $\beta_{60}=\SI{0.939}{\per\metre}$ --- from
a mode whose time constant $1/(\alpha_e\beta_1^2)$ is of order two centuries
down to one of order \SI{300}{\hour}.

\subsubsection{Transform pair and the norm}
The $\psi_n$ are orthogonal under the weight $r$ inherited from the cylindrical
Laplacian. Defining the norm
\begin{equation}
\mathcal{N}_n=\int_{r_{ce}}^{R_{\mathrm{far}}} r\,\psi_n^2(r)\,\mathrm{d}r ,
\label{eq:dbhe:norm}
\end{equation}
the transform pair is
\begin{equation}
\bar\theta_n(t)=\int_{r_{ce}}^{R_{\mathrm{far}}} r\,\psi_n(r)\,\theta(r,t)\,\mathrm{d}r
\qquad\text{(transform)}, \qquad
\theta(r,t)=\sum_{n=1}^{N}\frac{\psi_n(r)}{\mathcal{N}_n}\,\bar\theta_n(t)
\qquad\text{(inverse)} .
\label{eq:dbhe:pair}
\end{equation}
Only the wall value is ever needed downstream, so the code stores the
\emph{inversion weights} $\Omega_n=\psi_n(r_{ce})/\mathcal{N}_n$ once and
reconstructs the wall temperature as a single dot product.

The norms are evaluated by the trapezoidal rule on a radial grid that adapts to
the mode being integrated: at least $6000$ points, and at least $20$ points per
half-wavelength $\pi/\beta_n$. Resolving each oscillation is the whole issue ---
the integrand $r\psi_n^2$ oscillates $n$ times across the domain and is weighted
towards large $r$.

\paragraph{The non-positive-norm guard.}
Mathematically $\mathcal{N}_n>0$ always: it is the integral of a non-negative
quantity. The code nevertheless checks it and warns, naming the offending mode
indices, whenever $\mathcal{N}_n\le0$. The paranoia is not about the mathematics
but about the quadrature: a non-positive norm can only mean the radial grid has
aliased $\psi_n$ and returned a meaningless number. Because $\mathcal{N}_n$ sits
in the \emph{denominator} of every inversion weight, such a failure does not
produce a mildly wrong answer but a diverging one, so the warning names the two
correct remedies --- raise the grid resolution, or reduce the number of modes to
the band the grid can support. Silence here would let a discretisation failure
masquerade as a physical result.

\subsubsection{The transformed equations}
Multiply Eq.~\eqref{eq:dbhe:radial} by $r\psi_n$, integrate over
$[r_{ce},R_{\mathrm{far}}]$, and integrate the radial operator by parts twice.
Orthogonality annihilates every cross term and leaves
\begin{equation}
\frac{\mathrm{d}\bar\theta_n}{\mathrm{d}t} =-\alpha_e\beta_n^{2}\,\bar\theta_n
+\alpha_e\Big[r\big(\psi_n\theta'-\psi_n'\theta\big)\Big]_{r_{ce}}^{R_{\mathrm{far}}} .
\label{eq:dbhe:proj}
\end{equation}
The boundary term is where the physics enters, and both of our boundary
conditions were chosen so that it collapses. At $R_{\mathrm{far}}$ both
$\psi_n$ and $\theta$ vanish, so the upper limit contributes nothing. At
$r_{ce}$, $\psi_n'(r_{ce})=0$ kills the second piece, and Fourier's law
$-k_e\theta'(r_{ce})=q''=q'/(2\pi r_{ce})$ converts the first, leaving
\begin{equation}
\boxed{\; \frac{\mathrm{d}\bar\theta_n}{\mathrm{d}t}+\alpha_e\beta_n^{2}\,\bar\theta_n
=\alpha_e\beta_n^{2}\,F_n\,q' , \qquad F_n=\frac{\psi_n(r_{ce})}{2\pi k_e\,\beta_n^{2}} \; }
\label{eq:dbhe:modeODE}
\end{equation}
where $q'$ is the heat rate per unit borehole length delivered \emph{into} the
rock at the wall (so $q'<0$ during extraction) and $F_n$ is the code's
precomputed forcing coefficient. Writing the right-hand side as
$\alpha_e\beta_n^2 F_n q'$ rather than as a bare $F_n q'$ is deliberate: it makes
$F_n$ the \emph{steady gain} of the mode, $\bar\theta_n^\infty=F_n q'$, which is
what the tail correction of Section~\ref{sec:dbhe:tail} needs.

This decoupling is the entire point of the transform. A partial differential
equation in $(r,t)$ has become $N$ scalar ordinary differential equations that
do not talk to each other.

%---------------------------------------------------------------------
\subsection{Exact time march of the transformed modes}
\label{sec:dbhe:march}

Over one step $\Delta t$ the flux is held piecewise constant at its value
$q'^{\,k}$ from the fluid solve. Equation~\eqref{eq:dbhe:modeODE} is then a
linear first-order ODE with constant coefficients and integrates \emph{exactly}:
\begin{equation}
\boxed{\; \bar\theta_n^{\,k+1}=\bar\theta_n^{\,k}\,d_n +F_n\left(1-d_n\right)q'^{\,k},
\qquad d_n=\exp\!\left(-\alpha_e\beta_n^{2}\Delta t\right) \; }
\label{eq:dbhe:Urec}
\end{equation}
with $\Delta t=\SI{0.25}{\day}=\SI{6}{\hour}$ in the baseline. There is no
time-step stability limit and no truncation error in time: whatever error exists
comes from holding $q'$ constant across the step, not from the integrator. The
decay factors $d_n$ are computed once at initialisation.

Two properties are worth naming. First, the cost per step is $O(N)$ per axial
cell and \emph{no history is stored}: the state is the $N$ numbers
$\bar\theta_n$, the same size after twenty years as after one step. This is the
genuine advantage of the transform over a convolution formulation (g-function or
infinite cylindrical source), whose history buffer grows with simulated time. It
is \emph{not} that the transform avoids iteration --- the coupled
fluid--formation problem is solved by fixed point either way. Second,
$d_n\to1$ for slow modes and $d_n\to0$ for fast ones, so
Eq.~\eqref{eq:dbhe:Urec} degrades gracefully: fast modes relax immediately to
their quasi-steady value $F_n q'$, slow modes accumulate.

%---------------------------------------------------------------------
\subsection{The analytical tail correction}
\label{sec:dbhe:tail}

\subsubsection{Why a truncated series is not merely ``a bit less accurate''}
The high-order eigenfunctions are the ones with short radial wavelength, and
short radial wavelength is what resolves the steep temperature gradient hugging
the borehole wall at early time. Truncating at $N=60$ therefore does not blur
the solution uniformly; it removes the near-wellbore structure specifically. The
symptom is systematic: the retained modes reproduce too \emph{little} resistance
between wall and far field, so the model under-predicts how far the wall
temperature must drop to drive a given flux, and gets the first weeks to months
badly wrong.

\subsubsection{Deriving the deficit}
The deficit can be computed exactly, because both the true and the truncated
steady limits are known in closed form.

The true steady resistance is Eq.~\eqref{eq:dbhe:Rss}, the classical logarithmic
result for flux at $r_{ce}$ and fixed temperature at $R_{\mathrm{far}}$. What do
the $N$ retained modes give in the same limit? Setting
$\mathrm{d}\bar\theta_n/\mathrm{d}t=0$ in Eq.~\eqref{eq:dbhe:modeODE} gives
$\bar\theta_n^{\infty}=F_n q'$; inverting at the wall with
Eq.~\eqref{eq:dbhe:pair} and dividing by $q'$ gives the partial resistance
\begin{equation}
R_{\mathrm{partial}} =\frac{1}{q'}\sum_{n=1}^{N}\frac{\psi_n(r_{ce})}{\mathcal{N}_n}F_n q'
=\frac{1}{2\pi k_e}\sum_{n=1}^{N} \frac{\psi_n^{2}(r_{ce})}{\beta_n^{2}\,\mathcal{N}_n} .
\label{eq:dbhe:Rpartial}
\end{equation}
The identity being exploited is that the \emph{complete} series reproduces the
logarithm exactly,
\begin{equation}
\sum_{n=1}^{\infty}\frac{\psi_n^{2}(r_{ce})}{\beta_n^{2}\,\mathcal{N}_n}
=\ln\!\left(R_{\mathrm{far}}/r_{ce}\right),
\label{eq:dbhe:sumid}
\end{equation}
so the truncation error is not estimated --- it is the closed-form remainder
\begin{equation}
\boxed{\; R_{\mathrm{tail}}=R_{ss}-R_{\mathrm{partial}}
=\frac{\ln(R_{\mathrm{far}}/r_{ce})}{2\pi k_e} -\frac{1}{2\pi
k_e}\sum_{n=1}^{N}\frac{\psi_n^{2}(r_{ce})}{\beta_n^{2}\mathcal{N}_n}\;}
\label{eq:dbhe:Rtail}
\end{equation}
This quantity is not a rounding error. Because the series converges only as
$1/n^{2}$ under constant-flux forcing, evaluating
Eqs.~\eqref{eq:dbhe:Rss} and~\eqref{eq:dbhe:Rpartial} at the baseline
configuration ($N=60$, $R_{\mathrm{far}}=\SI{200}{\metre}$,
$r_{ce}=\SI{0.10795}{\metre}$) leaves $R_{\mathrm{tail}}$ at roughly a third of
$R_{ss}$. A model that ignored it would be wrong about the borehole-wall
temperature by a first-order amount, not a refinement.

\subsubsection{Carrying the deficit as one extra state}
The missing resistance is given its own lumped mode, advanced with exactly the
same exponential recursion but with the decay rate of the first \emph{omitted}
eigenvalue. The code estimates that eigenvalue by the asymptotic spacing of the
Bessel roots, $\beta_{N+1}\approx\beta_N+\pi/R_{\mathrm{far}}$:
\begin{equation}
\bar\theta_{\mathrm{tail}}^{\,k+1} =\bar\theta_{\mathrm{tail}}^{\,k}\,d_{\mathrm{tail}}
+R_{\mathrm{tail}}\left(1-d_{\mathrm{tail}}\right)q'^{\,k}, \qquad
d_{\mathrm{tail}}=\exp\!\left[-\alpha_e\left(\beta_N+\pi/R_{\mathrm{far}}\right)^{2}\Delta
t\right].
\label{eq:dbhe:Utail}
\end{equation}
Note that $R_{\mathrm{tail}}$ plays the role of $F_n$ here: it \emph{is} the
steady gain of the lumped remainder, which is why Eq.~\eqref{eq:dbhe:modeODE}
was normalised as it was. Assigning the whole remainder a single time constant
is an approximation --- the omitted modes span a range of rates --- but it is the
conservative one, since the slowest omitted mode governs when the lumped state
has finished responding.

The borehole-wall temperature is then reconstructed cell by cell as
\begin{equation}
T_{bw}(z)=T_\infty(z)+\sum_{n=1}^{N}\bar\theta_n\,\Omega_n +\bar\theta_{\mathrm{tail}},
\qquad \Omega_n=\frac{\psi_n(r_{ce})}{\mathcal{N}_n} .
\label{eq:dbhe:Tbw}
\end{equation}

\subsubsection{The optional second-set correction}
The notebook also contains a finer variant of the same idea: instead of lumping
the remainder into one state, compute a second block of eigenvalues (modes
$N+1$ through $N_B$, with $N_B=1200$) explicitly, and represent their
contribution by the quasi-steady resistance
\begin{equation}
C_B=\frac{1}{2\pi k_e}\sum_{n=N+1}^{N_B}
\frac{\psi_n^{2}(r_{ce})}{\beta_n^{2}\,\mathcal{N}_n} .
\label{eq:dbhe:CB}
\end{equation}
This is \textbf{disabled in the baseline} (\texttt{USE\_QS\_CORRECTION = False},
$C_B=0$), so the results reported here use the single-state analytical tail of
Eq.~\eqref{eq:dbhe:Utail}. The machinery is retained because it provides the
natural convergence check on the tail approximation: as $N_B$ grows, $C_B$ must
approach $R_{\mathrm{tail}}$.

%---------------------------------------------------------------------
\subsection{Fluid: the coaxial boundary-value problem}
\label{sec:dbhe:fluid}

\subsubsection{Quasi-steady hypothesis}
The fluid is treated as steady \emph{within} each time step. The separation of
scales is wide: the water transits the borehole in minutes to hours, the time
step is \SI{6}{\hour}, and the formation responds over months to decades. The
cost is that transients \emph{inside} the well --- the thermal front travelling
down the annulus after a step change in injection temperature, and the inertia
of the water column and steel --- are not represented. The model must therefore
not be asked about the first few hours after a start-up or a flow change; it is
a tool for seasonal and multi-year behaviour, and honest only at that scale.

\subsubsection{The two coupled energy balances}
Two counter-current streams exchange heat with each other through the tubing
wall --- the parasitic ``short circuit'' by which the hot returning water reheats,
or is cooled by, the descending water --- and the annulus additionally exchanges
with the borehole wall. Steady balances on differential elements, nondimensionalised
with $\zeta=z/L$ and
\begin{equation}
\theta_j=\frac{T_j-T_{\mathrm{in}}}{\Delta T_{sc}}, \qquad \Delta
T_{sc}(t)=\max\!\left(\max_z T_{bw}(z)-T_{\mathrm{in}},\;\SI{1}{\kelvin}\right),
\label{eq:dbhe:scale}
\end{equation}
give
\begin{align}
\frac{\mathrm{d}\theta_a}{\mathrm{d}\zeta}
&=\mathrm{NTU}_{bw}\left(\theta_{bw}-\theta_a\right)
-\mathrm{NTU}_{at}\left(\theta_a-\theta_t\right),
\label{eq:dbhe:bvpa}\\[2pt]
\frac{\mathrm{d}\theta_t}{\mathrm{d}\zeta}
&=-\,\mathrm{NTU}_{at}\left(\theta_a-\theta_t\right).
\label{eq:dbhe:bvpt}
\end{align}
The minus sign on the tubing equation is the counter-flow: $\zeta$ increases
downward for both equations, but the tubing stream travels upward, so the
derivative along its own direction of travel is $-\mathrm{d}/\mathrm{d}\zeta$.

The scale $\Delta T_{sc}$ is \emph{recomputed every step} from the current wall
field, not fixed once. The \SI{1}{\kelvin} floor prevents division by zero when
the well has been driven to thermal equilibrium with the injection temperature.

\subsubsection{The NTU groupings}
Both groups are conductance per unit length times length divided by stream
capacity rate, evaluated at the \emph{well} flow $\dot m_{\mathrm{well}}$ (not
the loop flow --- see Section~\ref{sec:dbhe:reynolds}):
\begin{equation}
\mathrm{NTU}_{at}=\frac{G_{a,t}^{\mathrm{eff}}\,L}{\dot m_{\mathrm{well}}c_{p,w}}, \qquad
\mathrm{NTU}_{bw}=\frac{G_{a,ce}^{\mathrm{eff}}\,L}{\dot m_{\mathrm{well}}c_{p,w}} .
\label{eq:dbhe:NTUdef}
\end{equation}
Each conductance is a series chain of resistances per unit length, built on
log-mean (thermodynamic centroid) radii
$\bar r=(r_o-r_i)/\ln(r_o/r_i)$ so that each solid wall can be split into an
inner and an outer half without introducing an error:
\begin{align}
\frac{1}{G_{a,t}^{\mathrm{eff}}} &=\underbrace{\frac{1}{2\pi
r_{t,o}h_a}+\frac{\ln(r_{t,o}/\bar r_t)}{2\pi k_t}}_{\text{annulus film + outer half of
tubing}} +\underbrace{\frac{\ln(\bar r_t/r_{t,i})}{2\pi k_t}+\frac{1}{2\pi
r_{t,i}h_{tb}}}_{\text{inner half of tubing + bore film}},
\label{eq:dbhe:Gat}\\[4pt]
\frac{1}{G_{a,ce}^{\mathrm{eff}}} &=\underbrace{\frac{1}{2\pi r_{ca,i}h_a}+\frac{\ln(\bar
r_{ca}/r_{ca,i})}{2\pi k_{ca}}}_{\text{annulus film + inner half of casing}}
+\underbrace{\frac{\ln(r_{ca,o}/\bar r_{ca})}{2\pi k_{ca}}+\frac{\ln(\bar
r_{ce}/r_{ca,o})}{2\pi k_{ce}}}_{\text{outer half of casing + inner half of cement}}
+\underbrace{\frac{\ln(r_{ce}/\bar r_{ce})}{2\pi k_{ce}}}_{\text{outer half of cement}} .
\label{eq:dbhe:Gace}
\end{align}
The steel conductivities ($k_{ca}=k_{t,\mathrm{bare}}=\SI{45}{\watt\per\metre\per\kelvin}$)
make the wall terms negligible in Eq.~\eqref{eq:dbhe:Gat} for a \emph{bare}
tubing, which is exactly the problem: the short circuit is then strong. Vacuum
insulated tubing, $k_t=\SI{0.02}{\watt\per\metre\per\kelvin}$, raises the tubing
resistance by more than three orders of magnitude and is what suppresses it. The
model carries both cases explicitly and labels the state ``Bare'' or
``Insulated'' from the value of $k_t$.

Water properties ($\rho$, $c_p$, $\mu$, $k_f$, $\mathrm{Pr}$) come from
IAPWS-IF97 at the operating pressure $P_{\mathrm{op}}=\SI{50}{bar}$, evaluated
once per well state at a representative mean temperature. A constant-property
shortcut was explicitly rejected in the code: a fixed \SI{70}{\celsius} estimate
misstates $\mu$ and $\mathrm{Pr}$ by about $-38\%$ at \SI{43}{\celsius} and
$+90\%$ at \SI{130}{\celsius}, and both feed the film coefficients directly.

%---------------------------------------------------------------------
\subsection{Transfer-matrix solution of the two-point BVP}
\label{sec:dbhe:transfer}

Equations~\eqref{eq:dbhe:bvpa}--\eqref{eq:dbhe:bvpt} are linear with
coefficients that are constant within each depth cell (the wall temperature
$\theta_{bw}$ is piecewise constant, one value per cell). Writing
$\boldsymbol{\theta}=(\theta_a,\theta_t)^{\mathsf T}$,
\begin{equation}
\frac{\mathrm{d}\boldsymbol{\theta}}{\mathrm{d}\zeta}
=\mathbf{A}\boldsymbol{\theta}+\mathbf{b}_k, \qquad \mathbf{A}=\begin{pmatrix}
-\left(\mathrm{NTU}_{at}+\mathrm{NTU}_{bw}\right) & +\mathrm{NTU}_{at}\\[2pt]
-\mathrm{NTU}_{at} & +\mathrm{NTU}_{at} \end{pmatrix}, \qquad
\mathbf{b}_k=\begin{pmatrix}\mathrm{NTU}_{bw}\,\theta_{bw,k}\\[2pt]
0\end{pmatrix},
\label{eq:dbhe:Amatrix}
\end{equation}
and the exact propagator across one cell of thickness $\Delta\zeta=1/N_Z$ is
\begin{equation}
\boldsymbol{\theta}_{k+1}=\boldsymbol{\Phi}\,\boldsymbol{\theta}_{k}
+\boldsymbol{\Psi}\,\mathbf{b}_k, \qquad \boldsymbol{\Phi}=e^{\mathbf{A}\Delta\zeta}, \qquad
\boldsymbol{\Psi}=\int_0^{\Delta\zeta}\!e^{\mathbf{A}s}\,\mathrm{d}s
=\mathbf{A}^{-1}\left(\boldsymbol{\Phi}-\mathbf{I}\right).
\label{eq:dbhe:PhiPsi}
\end{equation}
Both $\boldsymbol{\Phi}$ and $\boldsymbol{\Psi}$ are built once per flow state
and reused down the whole well. Because $\det\mathbf{A}=-\mathrm{NTU}_{at}\,
\mathrm{NTU}_{bw}$, the closed form for $\boldsymbol{\Psi}$ requires both
groups strictly positive; a perfectly insulating tubing ($k_t\to0$, so
$\mathrm{NTU}_{at}\to0$) would make $\mathbf{A}$ singular. The vacuum-insulated
case is small but finite, so this is a limit to be aware of rather than a
present defect.

\paragraph{Closing the boundary conditions.}
The problem is a two-point BVP: the annulus condition is known at the surface
($\theta_a(0)=0$, injection at $T_{\mathrm{in}}$) but the tubing condition is
known only at the bottom (the two streams join, $\theta_a(1)=\theta_t(1)$). The
code exploits linearity instead of shooting iteratively. It marches two
independent solutions from the surface,
\begin{equation}
\mathbf{y}^{P}_0=\begin{pmatrix}0\\
0\end{pmatrix} \;\text{(particular, with the }\mathbf{b}_k\text{ forcing)}, \qquad
\mathbf{y}^{H}_0=\begin{pmatrix}0\\
1\end{pmatrix} \;\text{(homogeneous, unforced)},
\label{eq:dbhe:seeds}
\end{equation}
both of which already satisfy $\theta_a(0)=0$, and forms
$\boldsymbol{\theta}=s\,\mathbf{y}^{H}+\mathbf{y}^{P}$. The single unknown $s$
is the surface tubing temperature, and the bottom joining condition fixes it in
one division:
\begin{equation}
s=\frac{y^{P}_{t}(1)-y^{P}_{a}(1)}{y^{H}_{a}(1)-y^{H}_{t}(1)} .
\label{eq:dbhe:shoot}
\end{equation}
No iteration, no stability limit, no step-size restriction --- the axial
discretisation error comes only from treating $\theta_{bw}$ as piecewise
constant, since the ODE itself is integrated exactly within each cell. The
produced temperature is read off directly:
\begin{equation}
T_{\mathrm{ret}}=T_{\mathrm{in}}+\theta_t(0)\,\Delta T_{sc}.
\label{eq:dbhe:Tret}
\end{equation}

%---------------------------------------------------------------------
\subsection{Closing the loop: flux back to the formation}
\label{sec:dbhe:coupling}

The fluid solve returns the annulus temperature at each cell centre, obtained by
averaging the two cell-face values,
$T_{a,cc,k}=T_{\mathrm{in}}+\tfrac12(\theta_{a,k}+\theta_{a,k+1})\Delta T_{sc}$.
That temperature sets the wall flux that drives Eq.~\eqref{eq:dbhe:Urec}:
\begin{equation}
\boxed{\; q'_k=G_{\mathrm{eff}}\left(T_{a,cc,k}-T_{bw,k}\right), \qquad
G_{\mathrm{eff}}=\frac{\mathrm{NTU}_{bw}\,\dot m_{\mathrm{well}}\,c_{p,w}}{L}
\;=\;G_{a,ce}^{\mathrm{eff}} \; }
\label{eq:dbhe:qeff}
\end{equation}
per unit borehole length. The last equality is a consistency identity, not a new
model: substituting Eq.~\eqref{eq:dbhe:NTUdef} recovers the conductance chain
of Eq.~\eqref{eq:dbhe:Gace} exactly. The same conductance therefore appears on
both sides of the coupling, which is what guarantees the energy the fluid gains
is the energy the rock loses. During extraction $T_{a,cc}<T_{bw}$, so $q'<0$,
and Eq.~\eqref{eq:dbhe:Urec} drives $\bar\theta_n$ negative and $T_{bw}$ below
$T_\infty$: the formation depletes.

The order within a step is \emph{evaluate, then advance}: the fluid is solved
against the wall field left by the previous step, and only after the surface
loop has converged is the formation marched with the resulting flux. Relative to
the formation's own time constants a \SI{6}{\hour} lag is negligible --- far
smaller than what the quasi-steady fluid hypothesis already concedes.

\paragraph{Retirement of the operating fraction $f_{\mathrm{op}}$.}
Earlier versions derated the forcing by a duty fraction,
$q'=f_{\mathrm{op}}\,G_{\mathrm{eff}}(T_{a,cc}-T_{bw})$, to represent a load
running only part of the day. That factor has been \textbf{removed} from
Eq.~\eqref{eq:dbhe:qeff}. In the coupled model the subsurface runs
\emph{continuously} and the duty cycle is expressed where it physically lives:
in the surface valve schedule, principally the well-flow fraction
$w_{v\mathrm{E}}$, and in the thermal accumulator. Derating the flux conflated
two distinct questions --- how hard the well is run, and how much load is served
--- which made the predicted formation depletion genuinely ambiguous. The
consequence must be stated plainly: the formation now sees the full, un-derated
flux, so for a given well flow the drawdown is faster than under the derated
formulation, and any earlier write-up carrying $f_{\mathrm{op}}$ inside the
transform forcing requires amendment. The symbol survives in the configuration
at $f_{\mathrm{op}}=1.00$ only so that legacy references do not break; it
multiplies nothing.

%---------------------------------------------------------------------
\subsection{Film coefficients, the Reynolds clamp, and the validity flag}
\label{sec:dbhe:reynolds}

Both film coefficients $h_a$ and $h_{tb}$ in
Eqs.~\eqref{eq:dbhe:Gat}--\eqref{eq:dbhe:Gace} come from the Gnielinski
correlation with the Petukhov friction factor,
\begin{equation}
f_D=\left(0.790\ln\mathrm{Re}-1.64\right)^{-2}, \qquad
\mathrm{Nu}=\frac{(f_D/8)\left(\mathrm{Re}-1000\right)\mathrm{Pr}}
{1+12.7\sqrt{f_D/8}\left(\mathrm{Pr}^{2/3}-1\right)}, \qquad h=\frac{\mathrm{Nu}\,k_f}{D_h},
\label{eq:dbhe:gnielinski}
\end{equation}
evaluated on the hydraulic diameters $D_{h,a}=2(r_{ca,i}-r_{t,o})$ for the
annulus and $D_{h,tb}=2r_{t,i}$ for the tubing bore. The two velocities are tied
by mass conservation, $v_{tb}=v_a A_a/A_{tb}$.

Equation~\eqref{eq:dbhe:gnielinski} is valid for $\mathrm{Re}\gtrsim3000$; below
that it is not merely inaccurate but pathological, since the factor
$(\mathrm{Re}-1000)$ drives $\mathrm{Nu}$ to zero at $\mathrm{Re}=1000$ and
negative below it. The well flow is $w_{v\mathrm{E}}\dot m$, and
$w_{v\mathrm{E}}$ is a control that the surface plant may throttle, so low
Reynolds numbers are reachable in normal operation. The code handles this in two
separate places, and the distinction matters.

\begin{enumerate}[leftmargin=*]
\item \emph{Numerically}, the correlation is \textbf{clamped}: the Reynolds
number fed to Eq.~\eqref{eq:dbhe:gnielinski} is $\max(\mathrm{Re},3100)$. Below
the clamp the film coefficient simply freezes at its value at
$\mathrm{Re}=3100$. This keeps $h$ finite and positive so the conductance chain
never produces a negative or singular $\mathrm{NTU}$, and it deliberately does
\emph{not} switch to a laminar correlation mid-march.
\item \emph{Diagnostically}, the true unclamped Reynolds numbers of both streams
are recomputed every time the flow is set, and a one-shot message is printed
naming $w_{v\mathrm{E}}$ and both values whenever either falls below $3000$. The
flag fires once per well state, so a long throttled run reports the condition
without flooding the log.
\end{enumerate}

Two facts make the clamp acceptable rather than fatal. First, the regime where
the correlation is weakest is the regime where the answer barely depends on it:
at low $w_{v\mathrm{E}}$ the capacity rate $\dot m_{\mathrm{well}}c_{p,w}$ is
small, $\mathrm{NTU}_{bw}$ is correspondingly large, the annulus water nearly
equilibrates with the wall regardless of $h_a$, and $T_{\mathrm{ret}}$ becomes
insensitive to it. Second, the control carries a floor
$w_{v\mathrm{E}}\ge w_{v\mathrm{E},\min}=0.05$, keeping a trickle circulating
rather than shutting down; a true shutdown would require a different physical
model (conduction into a stagnant column), and switching models mid-march would
introduce a discontinuity the tail correction of
Section~\ref{sec:dbhe:tail} has not been validated against. One continuous model
with an honest warning is the more defensible choice --- but the warning must be
read, because a run spending most of its time below $\mathrm{Re}=3000$ is
reporting film coefficients the correlation does not support.

%---------------------------------------------------------------------
\subsection{Summary of the subsurface parameters}
\label{sec:dbhe:params}

\begin{table}[htbp]
\centering
\caption{Baseline subsurface configuration, as declared in the model.}
\label{tab:dbhe:params}
\begin{tabular}{@{}llr@{}}
\toprule
Symbol & Quantity & Value \\
\midrule
\multicolumn{3}{@{}l}{\emph{Geometry (legacy S.\ Louisiana well: 2-7/8\,in tubing, 7\,in casing, 8.5\,in hole)}}\\
$r_{t,i}$      & tubing inner radius            & \SI{0.03099}{\metre}\\
$r_{t,o}$      & tubing outer radius            & \SI{0.03652}{\metre}\\
$r_{ca,i}$     & casing inner radius            & \SI{0.07975}{\metre}\\
$r_{ca,o}$     & casing outer radius            & \SI{0.08890}{\metre}\\
$r_{ce}$       & cement outer radius (hole)     & \SI{0.10795}{\metre}\\
\addlinespace
\multicolumn{3}{@{}l}{\emph{Materials and formation}}\\
$k_{ca}$       & casing steel conductivity      & \SI{45.0}{\watt\per\metre\per\kelvin}\\
$k_{ce}$       & cement conductivity            & \SI{0.95}{\watt\per\metre\per\kelvin}\\
$k_t$          & tubing: bare steel / VIT       & \num{45.0} / \SI{0.02}{\watt\per\metre\per\kelvin}\\
$k_e$          & formation conductivity         & \SI{2.2}{\watt\per\metre\per\kelvin}\\
$\alpha_e$     & formation diffusivity          & \SI{1.05e-6}{\metre\squared\per\second}\\
$T_{\mathrm{surf}}$ & mean surface temperature  & \SI{20.0}{\celsius}\\
$G_g$          & geothermal gradient            & \SI{0.031}{\kelvin\per\metre}\\
\addlinespace
\multicolumn{3}{@{}l}{\emph{Operation and numerics}}\\
$T_{\mathrm{in}}$ & injection temperature       & \SI{25.0}{\celsius}\\
$P_{\mathrm{op}}$ & wellbore operating pressure & \SI{50}{bar}\\
$L$            & depth (demonstration case)     & \SI{4500}{\metre}\\
$N_Z$          & axial cells in the fluid BVP   & 50\\
$N$            & retained eigenvalues (Set A)   & 60\\
$\Delta t$     & time step                      & \SI{0.25}{\day} = \SI{6}{\hour}\\
$R_{\mathrm{far}}$ & transform truncation radius & \SI{200}{\metre}\\
$w_{v\mathrm{E},\min}$ & floor on well-flow fraction & 0.05\\
$T_{\mathrm{cement}}^{\max}$ & cement reliability limit & \SI{165.0}{\celsius}\\
\bottomrule
\end{tabular}
\end{table}

Two derived quantities close the section, because both are load-bearing and
neither is a free parameter. The truncation radius follows from the horizon by
Eq.~\eqref{eq:dbhe:Rfar} and lands on its own floor at \SI{200}{\metre}; the
number of time steps follows from the horizon and $\Delta t$ and is
$\num{29220}$ for the twenty-year baseline. Every other number in
Table~\ref{tab:dbhe:params} is a declared input, and the two outputs the rest of
the plant consumes are Eq.~\eqref{eq:dbhe:Tret} and Eq.~\eqref{eq:dbhe:qeff}.
